\section{Discussion, Future Directions, and Conclusion}

In this paper, we proposed \MMCC{}, which gives tight amplification guarantees for sampling in the limit as $\epsilon \rightarrow 0$. One limitation of our work is that we are not able to prove adaptivity for non-lower triangular $\bfC$, which captures important matrix mechanisms like the ``fully efficient'' binary tree mechanism \citep{honaker15}. It is an important future direction to fully understand what combinations of privacy amplification and correlated noise allow the same privacy for non-adaptive and adaptive inputs. In addition, there are many potential improvements to \MMCC{}, as well as open problems that naturally follow from our work. 

First, our tail bound on the conditional sampling probabilities $\tilde{p}_{i,j}$ approach $p$ as $\sigma \rightarrow \infty$. However, for finite $\sigma$, $\tilde{p}_{i,j}$ can be much larger than $p$, i.e. the $\epsilon$ computed by \MMCC{} can be much larger than the true $\epsilon$. We believe the values of $\tilde{p}_{i,j}$ we compute are not tight and can be improved. In particular, in computing $\tilde{p}_{i,j}$, we give a tail bound on the maximum of the dot product of a Gaussian with a set of vectors; the values of $\tilde{p}_{i,j}$ we compute effectively correspond to the case where this tail bound is attained by every dot product. This is overly pessimistic; it should be possible to obtain tighter $\epsilon$ via a more refined tail-bounding approach. 

Second, while \MMCC{} has a polynomial dependence on $n$ (whereas computing $H_\epsilon$ via e.g. numerical integration would require time exponential in $n$), empirically we found that even with many optimizations for runtime, running \MMCC{} for $n \approx 2000$ still took several hours. In practice, we would often like to run for larger $n$, or do multiple sequential runs of \MMCC{} in order to e.g. compute the smallest $\sigma$ that gives a certain $\epsilon$ via binary search. In turn, it is practically interesting and important to make \MMCC{} more computationally efficient or discover alternative algorithms that give comparable $\epsilon$ at smaller runtime.

Our interest in the matrix mechanism is primarily motivated by the works of \cite{denisovDPMF, MEMF, BandedMF} which considered the problem of choosing $\bfC$ that optimizes (a proxy for) the utility of DP-FTRL. The utility of DP-FTRL can be written as a function of $\bfC^{-1}$, and thus can be optimized under a constraint of the form ``the matrix mechanism defined by $\bfC$ satisfies a given privacy definition''. Without amplification, this constraint can usually be easily written as e.g. $\bfC \in \calS$ where $\calS$ is a convex set of matrices, which makes optimizing under this constraint easy. An interesting question is whether we can solve the same problem but accounting for amplification. This would likely require designing a function that takes $\bfC, p, \sigma$ and approximates $\epsilon$ that is differentiable in $\bfC$ (unlike \MMCC{} because it is an algorithmic computation that is not easily differentiable).

In these works, DP-FTRL is always strictly better than DP-SGD without amplification, but with amplification for small $\epsilon$ the optimal choice of $\bfC$ with amplification is the identity, i.e. the optimal DP-FTRL is just DP-SGD (with independent noise). If we could optimize $\bfC$ under an amplified privacy constraint, we conjecture the following (perhaps surprising) statement could be proven as a corollary: As long as we are not in the full-batch setting, even with amplification by sampling, the optimal choice of $\bfC$ is never the identity for $\epsilon > 0$. In other words, despite its ubiquity, DP-SGD is never the optimal algorithm to use (ignoring computational concerns).